\section{Análisis del enfoque}
\label{cap:approachAnalysis}

\subsection{Representación Vectorial Contextual en Espacio Compartido}
\label{sec:nlp_representation}

Para vectorizar los textos del corpus sin sufrir las limitaciones de los modelos de bolsa de palabras (TF-IDF) ni la falta de contexto de los \textit{word embeddings} estáticos, se seleccionan modelos de lenguaje basados en transformadores multilingües (\textit{Sentence-Transformers} como \textit{paraphrase-multilingual-mpnet-base-v2} o \textit{multilingual-e5-large}). Estos modelos permiten proyectar los títulos y descriptores modulares de las asignaturas de las mallas, las oraciones de los perfiles de egreso declarados, las definiciones de las \textit{Knowledge Areas} de la ACM/IEEE y las secciones de requerimientos de las vacantes de LinkedIn dentro de un \textbf{mismo espacio vectorial latente continuo ($\mathbb{R}^d$)}. Esta proyección homotópica habilita el cálculo directo de distancia euclidiana y similitud de coseno ($\cos \theta$) entre elementos de distintas fuentes de datos en español e inglés.

La alineación y las brechas semánticas entre una malla curricular $M_i$ y una fuente de referencia (ACM o LinkedIn) se cuantifican mediante tres métricas principales:
\begin{enumerate}[label=\alph*)]
	\item \textbf{Similitud de Coseno ($\cos \theta$):} Mide la orientación relativa entre el vector de una asignatura $\mathbf{v}_{\text{asig}}$ y el vector de una competencia o área de conocimiento $\mathbf{v}_{\text{ref}}$:
	      \begin{equation}
		      \cos(\mathbf{v}_{\text{asig}}, \mathbf{v}_{\text{ref}}) = \frac{\mathbf{v}_{\text{asig}} \cdot \mathbf{v}_{\text{ref}}}{\|\mathbf{v}_{\text{asig}}\| \|\mathbf{v}_{\text{ref}}\|}
	      \end{equation}
	\item \textbf{Tasa de Cobertura de la Malla ($\% \text{Cobertura}$):} Representa el porcentaje de conceptos clave o \textit{Knowledge Areas} del estándar internacional o del mercado laboral representados en la malla curricular con una similitud semántica $\cos \theta \ge 0.60$.
	\item \textbf{Índice de Brecha Semántica ($1 - \cos \theta$):} Cuantifica la distancia vectorial que refleja el déficit o la ausencia explícita de asignaturas en la malla curricular para cubrir requerimientos tecnológicos específicos.
\end{enumerate}

\subsection{Justificación de la Estrategia Tripartita de Agrupamiento}
\label{sec:clustering_justification}

Para descubrir la estructura disciplinar de las mallas curriculares y agrupar las asignaturas de forma no supervisada, se adopta una \textbf{estrategia comparativa tripartita} que combina K-Means, Agrupamiento Jerárquico (\textit{Hierarchical Clustering}) y DBSCAN. Esta combinación no es arbitraria; responde a que cada algoritmo resuelve una necesidad analítica complementaria e insustituible sobre los datos curriculares:
\begin{enumerate}[label=\arabic*.]
	\item \textbf{K-Means Clustering:} Produce particiones convexas en $k$ grupos maximizando la homogeneidad interna. Su principal ventaja matemática radica en que calcula centroides explícitos $\mathbf{c}_k \in \mathbb{R}^d$ para cada clúster, los cuales actúan como la representación vectorial promedio de un área disciplinar dentro de la malla. El número óptimo de clústeres $k$ se determina mediante la maximización del Coeficiente de Silhouette.
	\item \textbf{Agrupamiento Jerárquico Aglomerativo (Ward Linkage):} Construye una estructura arborescente o dendrograma que preserva las relaciones de jerarquía y subordinación taxonómica entre asignaturas (e.g., relacionando asignaturas de programación básica con materias de ingeniería de software avanzada). Permite visualizar los niveles de abstracción disciplinar sin prefijar el número de grupos.
	\item \textbf{DBSCAN (Density-Based Spatial Clustering of Applications with Noise):} Algoritmo basado en densidad que define clústeres como regiones densas de vectores en $\mathbb{R}^d$ separadas por zonas de baja densidad. Su función técnica crucial en este proyecto es la \textbf{detección y aislamiento explícito de asignaturas atípicas o ruido (\textit{outliers})}, es decir, asignaturas singulares presentes en mallas de ciertas universidades que no encajan en las agrupaciones estándar del sistema.
\end{enumerate}

Para evaluar la validez externa de los clústeres generados frente a un estándar objetivo, se realiza una anotación manual de las 747 asignaturas del corpus clasificándolas en las 18 \textit{Knowledge Areas} de la ACM (CS2013). Esta clasificación actúa como etiqueta de referencia (\textit{ground truth}) para calcular el Índice Rand Ajustado ($\text{ARI} \ge 0.35$), la Información Mutua Normalizada ($\text{NMI} \ge 0.40$) y la Pureza del agrupamiento ($\text{Purity} \ge 0.70$).
